% Options for packages loaded elsewhere
\PassOptionsToPackage{unicode}{hyperref}
\PassOptionsToPackage{hyphens}{url}
%
\documentclass[
]{article}
\usepackage{amsmath,amssymb}
\usepackage{lmodern}
\usepackage{iftex}
\ifPDFTeX
  \usepackage[T1]{fontenc}
  \usepackage[utf8]{inputenc}
  \usepackage{textcomp} % provide euro and other symbols
\else % if luatex or xetex
  \usepackage{unicode-math}
  \defaultfontfeatures{Scale=MatchLowercase}
  \defaultfontfeatures[\rmfamily]{Ligatures=TeX,Scale=1}
\fi
% Use upquote if available, for straight quotes in verbatim environments
\IfFileExists{upquote.sty}{\usepackage{upquote}}{}
\IfFileExists{microtype.sty}{% use microtype if available
  \usepackage[]{microtype}
  \UseMicrotypeSet[protrusion]{basicmath} % disable protrusion for tt fonts
}{}
\makeatletter
\@ifundefined{KOMAClassName}{% if non-KOMA class
  \IfFileExists{parskip.sty}{%
    \usepackage{parskip}
  }{% else
    \setlength{\parindent}{0pt}
    \setlength{\parskip}{6pt plus 2pt minus 1pt}}
}{% if KOMA class
  \KOMAoptions{parskip=half}}
\makeatother
\usepackage{xcolor}
\setlength{\emergencystretch}{3em} % prevent overfull lines
\providecommand{\tightlist}{%
  \setlength{\itemsep}{0pt}\setlength{\parskip}{0pt}}
\setcounter{secnumdepth}{-\maxdimen} % remove section numbering
\ifLuaTeX
  \usepackage{selnolig}  % disable illegal ligatures
\fi
\IfFileExists{bookmark.sty}{\usepackage{bookmark}}{\usepackage{hyperref}}
\IfFileExists{xurl.sty}{\usepackage{xurl}}{} % add URL line breaks if available
\urlstyle{same} % disable monospaced font for URLs
\hypersetup{
  hidelinks,
  pdfcreator={LaTeX via pandoc}}

\author{Dietmar Gerald Schrausser}
\date{09.05.2026}


\begin{document}

\hypertarget{foliumblat-vr}{%
\section{Foliu{[}m{]}/Blat Vr}\label{foliumblat-vr}}

Eine zeitgenössische deutsche Übersetzung wird präsentiert, die auf
einem Vergleich der bestehenden lateinischen, frühneuhochdeutschen
(ENHG) und englischen (Übersetzungs-) Texte basiert, um ein besseres
Verständnis zu ermöglichen. Insbesondere da der ENHG-Text im Vergleich
zum modernen Hochdeutschen schwer verständlich ist, da dieser fundierte
Deutschkenntnisse (als Muttersprache) sowie Kenntnisse verschiedener
Dialekte voraussetzt.

\hypertarget{vom-werk-des-sechsten-tages}{%
\subsection{Vom Werk des sechsten
Tages}\label{vom-werk-des-sechsten-tages}}

\begin{quote}
\textbf{S}Exte die dixit de{[}us{]}. p{[}ro{]}ducat terra
a{[}n{]}i{[}m{]}am viue{[}n{]}te{[}m{]}. iume{[}n{]}ta. {[}et{]}
reptila. {[}et{]} bestias terre iuxta spe{[}cie{]}s suas. Et vidit
de{[}us{]} q{[}uo{]}d e{[}ss{]}et bonu{[}m{]}: ait. Faciam{[}us{]}
ho{[}m{]}i{[}n{]}em ad imagine{[}m{]} {[}et{]} si{[}mi{]}litudine{[}m{]}
n{[}ost{]}ram. Et p{[}er{]}sit piscib{[}us{]} mar{[}is{]}. {[}et{]}
volatilib{[}us{]} celi. {[}et{]} bestijs vniuerse terre. Et creuit
de{[}us{]} ho{[}m{]}i{[}n{]}em ad imagine{[}m{]} {[}et{]}
si{[}mi{]}litudine{[}m{]} sua{[}m{]}.
\end{quote}

Nachdem die oberen Teile der Welt geschmückt waren, schmückte er
schließlich am sechsten Tag die Erde mit den Geschlechtern der
Lebewesen\footnote{`generib{[}us{]} a{[}n{]}i{[}m{]}aliu{[}m{]}' und
  `geschlechte{[}n{]} der thier'.}. Unter den \emph{Tieren} der Erde
erwähnt \textbf{Moses} drei: (i) \emph{Lasttiere}, (ii)
\emph{kriechende} und (iii) \emph{wilde Tiere}\footnote{`iume{[}n{]}ta.
  reptilia. {[}et{]} bestias' und `iohthier kriechende vnd wildthier'.}.
Er deutet damit allgemein auf drei verschiedene Arten von
\emph{unvernünftigen}\footnote{`irr{[}ati{]}onabiliu{[}m{]}' und
  `unvernüftigen'.} Tieren hin. \emph{Wildtiere} (iii) mit vollkommener
\emph{Vorstellungskraft}\footnote{`pha{[}n{]}tasia' und `fantasey', in
  der mittelalterlichen Philosophie bezeichnet dies die
  Vorstellungskraft.} nehmen unter den unvernünftigen
\emph{Dingen}\footnote{`sortite' und `thieren'.} eine Mittelstellung
ein, obgleich sie sich vom Menschen weder belehren noch zähmen lassen.
So stehen die \emph{kriechenden Tiere} (ii) mit \emph{un}vollkommener
\emph{Vorstellungskraft} zwischen dem Vieh\footnote{`bruta' und `vieh'.}
und Pflanzen. \emph{Lasttieren} (i) mangelt es zwar an Vernunft, können
aber dennoch vom Menschen weitgehend gezähmt werden, so scheinen sie an
\emph{etwas} Vernunft teilzuhaben. Diese nehmen gleichsam eine
Zwischenstellung zwischen Vieh\footnote{`bruta' und `vihe'.} und den
Menschen ein.

Weiters ordnete er an, dass \emph{Lebewesen}\footnote{`A{[}n{]}i{[}m{]}alia'
  und `thier'.} aller Art und Gestalt, groß und klein, paarweise
geschaffen werden sollten, von beiderlei Geschlecht\footnote{`diuersi
  sex{[}us{]} sing{[}u{]}la' und `bede me{[}n{]}dlein vnd freülein'.}.
Von ihren Nachkommen wurden Luft, Erde und Meere erfüllt. Gott gab allen
von Generation zu Generation Ernährung von der Erde, so könnten sie auch
den Menschen \emph{nützlich} sein: Überzählige(s) als
Nahrung\footnote{`nimi{[}orum{]} ad cibos' und `als etlich zu speysung',
  wohl die tierischen \emph{Produkte} wie Eier oder Milch.}, wahrlich
sonst für Kleidung\footnote{`v{[}er{]}o ad vestime{[}n{]}tu{[}m{]}' und
  `ettlich zebeklaidung', wie Wolle etc.}, diejenigen, die über große
Kraft verfügen, um beim Pflügen des Landes zu helfen, aus diesem Grund
werden diese auch (im Lateinischen) \emph{iumenta}\footnote{`vt i{[}n{]}
  excole{[}n{]}da terra iuuarent. vn{[}de{]} dicta sunt iume{[}n{]}ta.'}
genannt\footnote{aus den `Etymologiae' (De Seville,
  \href{https://books.google.com/books?id=wb8Z_vlSyJwC}{1802}, Buch XII,
  1.7).}.

Bislang haben wir (erst) von \emph{drei} Welten gesprochen, der (i)
\emph{überhimmlischen}, der (ii) \emph{himmlischen} und der (iii)
\emph{sublunaren}, nun wollen wir den (iv) \emph{Menschen} als Teil
einer vierten Welt betrachten. Nachdem nun alles in wunderbarer Weise
geordnet war, beschloss er sich ein ewiges Reich zu bereiten, im Voraus
unzählige Seelen zu erschaffen, denen er Unsterblichkeit verleihen
konnte. Daraufhin schuf er sich ein Gleichnis seiner Selbst, empfindsam
und intelligent, geformt nach seinem eigenen Bild\footnote{`imaginis sue
  forma{[}m{]}' und `zu form oder gestalt seiner pildnus'.}, dem nichts
vollkommener gleicht. So formte er den Menschen aus dem Lehm der Erde,
daher auch der Name \emph{Mensch}\footnote{`ho{[}mo{]}' und
  `mensche{[}n{]}'}, geschaffen aus \emph{Erde}\footnote{`humo' und
  `lette{[}n{]} oder kloße der erde{[}n{]}'.}.\footnote{spiegelt eine
  gängige mittelalterliche etymologische Vorstellung wider, die
  \emph{Homo} mit \emph{Humus} in Verbindung bringt.}

Gott, der Schöpfer aller Dinge, schuf den Menschen, (was auch
\textbf{Cicero} erkannte, obwohl er keine himmlischen Schriften besaß)
\emph{dieser} überlieferte im ersten Buch über die Gesetze dasselbe wie
auch die Propheten, seine Worte haben wir im Folgenden
wiedergegeben:\footnote{vgl. die `Divinae Institutiones' (Lactantius,
  \href{https://books.google.com/books?id=qGsnghsi3uUC}{1497}).}

\begin{quote}
Hoc a{[}n{]}i{[}m{]}al p{[}ro{]}uidu{[}m{]}: sagax: m{[}u{]}ltiplex:
acutu{[}m{]}: memor: plenu{[}m{]} r{[}ati{]}o{[}n{]}is {[}et{]}
{[}con{]}silij. que{[}m{]} vocam{[}us{]} homine{[}m{]}: p{[}re{]}clara
q{[}ua{]}da{[}m{]} {[}con{]}dit{[}i{]}o{[}n{]}e generatu{[}m{]}
e{[}ss{]}e a summo deo solu{[}m{]}.\\
Est e{[}ni{]}m ex tot a{[}n{]}i{[}m{]}antiu{[}m{]} generib{[}us{]}
at{[}que{]} natur{[}is{]}: p{[}ar{]}ticeps r{[}ati{]}o{[}n{]}is {[}et{]}
cogitat{[}i{]}o{[}n{]}is. cu{[}m{]} cetera sint o{[}mn{]}ia
exp{[}er{]}tia.\footnote{`Dieses Lebewesen(!) das wir Mensch nennen:
  vorsichtig, flink, vielfältig, scharf im Gedächtnis, voll Vernunft und
  Rat; in klarer Art und Eigenschaft allein von dem höchsten Gott
  geboren. Denn es hat an allen Geschlechtern und Naturen der beseelten
  Geschöpfe teil, {[}besonders aber{]} an der Vernunft und dem Verstand,
  was allen anderen Geschöpfen fehlt', `De Legibus' (Cicero \& Cicero,
  \href{https://doi.org/10.3931/e-rara-89116}{1496}, Buch I, Abschnitt
  22), bespricht den Menschen als in einzigartiger Weise mit Vernunft
  begabt, welche ihn mit dem \emph{Göttlichen} verbinde.}
\end{quote}

Übrigens ist ein weit verbreiteter Brauch unter Königen oder Fürsten der
Erde, dass sie, würden sie eine \emph{prächtige}, \emph{edle} Stadt
erdacht und gebaut haben, nach der Fertigstellung ihr eigenes Bildnis in
der Stadtmitte errichteten, dass es von allen gesehen und betrachtet
werden könne. Ebenso hat auch Gott, der \emph{Fürst} aller Dinge,
gehandelt, als er, nachdem das gesamte Gefüge\footnote{`machi{[}n{]}a'
  und `paw'.} der Welt konstruiert war, schließlich den Menschen
\emph{inmitten} von all diesem setzte, geschaffen nach seinem Bild und
in seinem Gleichnis geformt, dass man es wohl wie \emph{dieser}
\textbf{Merkur} ausrufen möchte:\footnote{vgl. `De Vita Coelitus
  Comparanda' (Ficinus,
  \href{https://doi.org/10.3931/e-rara-68641}{1497}).}

\begin{quote}
Magnu{[}m{]} o asclepi mirac{[}u{]}l{[}u{]}m e{[}st{]} ho{[}mo{]}.\\
Hoc p{[}re{]}cipuo no{[}m{]}i{[}n{]}e gl{[}or{]}iari hu{[}m{]}ana
{[}con{]}ditio p{[}otes{]}t: q{[}uo{]} et{[}iam{]}
f{[}a{]}ct{[}u{]}m{]}: vt s{[}er{]}uire illi: nulla creata
s{[}u{]}b{[}stanti{]}a dedignet{[}ur{]}: huic terra et elem{[}en{]}ta:
huic bruta sunt p{[}re{]}sto: {[}et{]} famulant{[}ur{]}. huic militat
celu{[}m{]}: huic salute{[}m{]}: p{[}ro{]}cura{[}n{]}t a{[}n{]}gelices
me{[}n{]}tes.\\
Nec miru{[}m{]} alicui videri debet: amari illu{[}m{]} ab
om{[}n{]}ibus.\\
In quo om{[}n{]}ia suu{[}m{]} aliq{[}ui{]}d:\\
Immo se tota {[}et{]} sua om{[}n{]}ia agnoscu{[}n{]}t.\footnote{`O
  Asklepios, der Mensch ist ein großes Wunder. Durch diesen
  außergewöhnlichen Namen kann sich das menschliche Dasein rühmen: Denn
  es ist geschehen, dass kein Geschöpf es verschmäht, ihm zu dienen.
  Erde und Elemente stehen ihm zur Seite; die Tiere sind bereit für ihn
  und dienen ihm. Für ihn kämpfen die Himmel; für ihn bringen Engel
  Heil. Es sollte niemanden verwundern, dass er von allen geliebt wird,
  in dem alles etwas Eigenes findet; ja, alles und alles Eigene erkennen
  sie in ihm', aus dem `Asklepios', einem hermetischen Text (vgl.
  Trismegistu,
  \href{https://gdz.sub.uni-goettingen.de/id/PPN858183994}{1471}), auch
  zitiert in Picos `Oratio de hominis dignitate' (Mirandola \&
  Mirandola, \href{https://doi.org/10.3931/e-rara-61217}{1601}).}
\end{quote}

\hypertarget{references}{%
\subsection{References}\label{references}}

Cicero, M. T., \& Cicero, Q. (1496). \emph{M.T. Ciceronis de Natura
Deorum Libri Tres} \ldots{} Impræssum Venetiis: per Symonem Papiensem
dictum Biuilaqua. \url{https://doi.org/10.3931/e-rara-89116}

De Seville, I. (1802). \emph{S. Isidori Hispalensis Episcopi Hispaniarum
Doctoris Opera Omnia}. Edited by Arevalo, F. Romae: Apvd Antonivm
Fvlgonivm. \url{https://books.google.com/books?id=wb8Z_vlSyJwC}

Ficinus, M. (1497). \emph{De Triplici Vita}. Basel: Johann Amerbach.
\url{https://doi.org/10.3931/e-rara-68641}

Lactantius, L.C.F. (1497). \emph{Divinae Institutiones}. Venetiis: Simon
Bevilaqua. \url{https://books.google.com/books?id=qGsnghsi3uUC}

Mirandola, G., \& Mirandola, G. F. (1601). \emph{Ioannis Pici,
Mirandulae \ldots{} opera quae extant omnia \ldots{}} Basileae: per
Sebastianum Henricpetri. \url{https://doi.org/10.3931/e-rara-61217}

Trismegistu, H. (1471). \emph{Mercurii Trismegisti Liber de Potestate Et
Sapientia Dei e Graeco in Latinum Traductus a Marsilio Ficino Pimander
Incipit}. Edited by Ficinus, M. Treviso.
\url{https://gdz.sub.uni-goettingen.de/id/PPN858183994}

\end{document}
